\chapter{AgentOS 运行时的信号本体框架}

\section{从模块设计转向信号设计}

传统软件系统通常首先从模块出发理解系统结构：

\[
System
=
Module_1
+
Module_2
+
\cdots
+
Module_n.
\]

这种方法适合描述静态的软件组成，却不足以刻画一个持续运行、持续感知、持续作用并持续改变自身结构的智能体。

对于持续智能体而言，更基础的问题不是：

\[
\text{系统包含哪些模块？}
\]

而是：

\[
\boxed{
\text{系统内部究竟有哪些不同性质的信号在流动？}
}
\]

因为任何模块最终都只是某些信号的：

\[
\text{接收者},
\qquad
\text{变换器},
\qquad
\text{产生者},
\qquad
\text{路由节点}.
\]

因此，在 AgentOS 的第一性设计中，应当首先建立信号本体，再讨论具体功能模块。

本章提出，AgentOS 运行时至少存在四类具有不同因果角色的基本信号：

\[
\boxed{
\mathcal S
=
\mathcal S_S
\oplus
\mathcal S_I
\oplus
\mathcal S_U
\oplus
\mathcal S_C
}
\]

分别为：

\[
\boxed{
\mathcal S_S:
\text{状态信号}
}
\]

\[
\boxed{
\mathcal S_I:
\text{信息信号}
}
\]

\[
\boxed{
\mathcal S_U:
\text{控制信号}
}
\]

\[
\boxed{
\mathcal S_C:
\text{信用信号}
}
\]

四者分别回答四个不同的问题：

\[
\boxed{
\text{状态：系统现在是什么？}
}
\]

\[
\boxed{
\text{信息：系统现在获得或表达了什么？}
}
\]

\[
\boxed{
\text{控制：系统接下来应该怎样变化？}
}
\]

\[
\boxed{
\text{信用：过去哪些结构应该因为结果而获得修改责任？}
}
\]

这四类信号共同构成 AgentOS 运行时最基本的信号本体。

\section{信号本体不是数据格式分类}

首先需要澄清：

\[
\boxed{
\text{信号类型}
\neq
\text{数据类型}.
}
\]

一个向量：

\[
x\in\mathbb R^d
\]

既可能表示状态，也可能表示信息，还可能表示控制量或信用量。

因此，不能通过其物理载体判断信号本体。

真正决定信号类别的是：

\[
\boxed{
\text{它在系统因果结构中承担什么作用。}
}
\]

例如同样一个标量：

\[
0.8
\]

如果表示当前资源占用比例，它是状态信号；

如果表示某个外部对象出现的概率，它是信息信号；

如果表示执行某个动作的控制增益，它是控制信号；

如果表示某个内部结构对一次失败应承担的责任比例，它是信用信号。

因此定义信号类别的不是：

\[
RepresentationFormat,
\]

而是：

\[
\boxed{
CausalRole.
}
\]

由此可以定义一个信号：

\[
\sigma
=
(x,\rho,\tau,\kappa),
\]

其中：

\[
x
\]

为具体载荷；

\[
\rho
\]

为因果角色；

\[
\tau
\]

为有效时间尺度；

\[
\kappa
\]

为作用范围。

信号本体主要由：

\[
\rho
\]

决定。

\section{状态信号：系统当前存在方式的表达}

状态信号描述系统在某个时刻实际处于什么状态。

设 AgentOS 的运行状态空间为：

\[
\mathcal X.
\]

则：

\[
x_t\in\mathcal X
\]

表示时刻 $t$ 的系统状态。

状态信号的核心定义为：

\[
\boxed{
\mathcal S_S
=
\text{对当前系统状态及其局部变化的可观测表达}.
}
\]

典型状态可以包括：

\[
x_t
=
[
x_t^{internal},
x_t^{resource},
x_t^{task},
x_t^{environment},
\ldots
].
\]

状态信号回答的不是：

\[
\text{“应该做什么？”}
\]

而是：

\[
\boxed{
\text{“现在是什么样？”}
}
\]

例如：

当前计算负载；

当前任务阶段；

当前资源占用；

当前内部活动强度；

当前环境状态；

当前执行进度。

都属于状态描述。

从动力系统角度，状态具有一个重要性质：

\[
\boxed{
x_t
}
\]

与输入和控制共同决定：

\[
x_{t+\Delta t}.
\]

因此系统最基础的状态方程可以写成：

\[
\boxed{
x_{t+\Delta t}
=
F(
x_t,
i_t,
u_t;
\Xi_t
).
}
\]

其中：

\[
i_t
\]

为信息信号；

\[
u_t
\]

为控制信号；

\[
\Xi_t
\]

表示当前系统中相对稳定、参与状态生成的内部结构。

所以状态不是静态数据库中的一条记录，而是：

\[
\boxed{
\text{系统当前动力学的承载面}.
}
\]

\section{信息信号：关于世界与系统的结构表达}

信息信号与状态信号不同。

状态信号描述：

\[
\text{系统当前是什么}.
\]

信息信号则描述：

\[
\text{系统当前获得、生成或调用了什么结构内容}.
\]

设信息表示空间为：

\[
\mathcal I.
\]

则：

\[
i_t\in\mathcal I.
\]

信息可能来自：

外部环境；

历史经验；

内部预测；

内部结构激活；

其他智能体；

工具或外部系统。

因此可以写成：

\[
i_t
=
[
i_t^{external},
i_t^{internal},
i_t^{predicted},
i_t^{retrieved},
\ldots
].
\]

信息信号的核心定义为：

\[
\boxed{
\mathcal S_I
=
\text{能够改变系统对可作用现实之内部表示的结构性信号}.
}
\]

这意味着：

\[
Information
\neq
RawData.
\]

原始输入只有在进入 Agent 的结构表示以后，才真正成为 AgentOS 意义上的信息。

可以定义一个一般观察映射：

\[
O_t:
\mathcal U_{accessible}
\rightarrow
\mathcal I,
\]

其中：

\[
\mathcal U_{accessible}
\]

表示 Agent 当前能够接触的现实部分。

因此：

\[
i_t
=
O_t(U_t).
\]

但：

\[
\boxed{
i_t\neq U_t.
}
\]

内部信息不是现实本身，而是现实经过：

\[
Interaction
\rightarrow
Projection
\rightarrow
Compression
\rightarrow
Representation
\]

之后形成的系统内部结构。

这一点对于整个 AgentOS 至关重要：

\[
\boxed{
Reality
\neq
InformationRepresentation.
}
\]

\section{状态与信息为什么必须分离}

状态和信息经常同时以向量形式存在，因此容易被混淆。

但是两者的因果角色不同。

考虑一个 Agent 正在处理一个高复杂度任务。

信息信号可能是：

\[
i_t=
\text{“当前输入包含一个复杂证明问题”}.
\]

而状态信号可能是：

\[
x_t=
\text{“当前内部资源已经接近承载上限”}.
\]

前者描述：

\[
\boxed{
\text{正在处理什么}.
}
\]

后者描述：

\[
\boxed{
\text{系统现在处于什么运行状态}.
}
\]

因此：

\[
\boxed{
Information
\neq
State.
}
\]

然而两者又持续耦合。

信息改变状态：

\[
i_t
\rightarrow
x_{t+\Delta t}.
\]

状态反过来决定信息如何被解释：

\[
x_t
\rightarrow
Interpret(i_t).
\]

所以更准确地写：

\[
\boxed{
x_{t+\Delta t}
=
F(
x_t,
I(i_t|x_t),
u_t
).
}
\]

即：

\[
\boxed{
\text{信息并不是脱离状态独立发挥作用，而是在状态条件下被解释。}
}
\]

\section{控制信号：改变系统状态转移方式}

第三类基本信号是控制信号。

设控制空间：

\[
\mathcal U.
\]

控制信号：

\[
u_t\in\mathcal U.
\]

其基本作用不是描述当前状态，也不是承载世界内容，而是：

\[
\boxed{
\text{改变系统接下来如何演化}.
}
\]

所以控制的本体定义可以写成：

\[
\boxed{
\mathcal S_U
=
\text{对状态转移过程施加方向性约束或调制的信号}.
}
\]

如果没有控制：

\[
x_{t+\Delta t}
=
F(x_t,i_t).
\]

加入控制以后：

\[
\boxed{
x_{t+\Delta t}
=
F(x_t,i_t,u_t).
}
\]

这里：

\[
u_t
\]

可以改变：

状态转移方向；

状态转移速度；

资源分配；

激活强度；

搜索范围；

信息进入比例；

内部过程优先级。

所以控制信号回答：

\[
\boxed{
\text{“系统下一步应该怎样运行？”}
}
\]

而不是：

\[
\text{“系统现在是什么？”}
\]

\section{控制信号不是信息信号}

一个命令本身可能具有语义内容，但其系统角色不只是提供信息。

例如：

\[
\text{“降低当前搜索深度”}
\]

如果只是作为一句文本被读取，它是信息。

但当其真正改变：

\[
F
\]

中的运行参数时，它就产生了控制作用。

因此可以定义从信息到控制的映射：

\[
\boxed{
\Gamma:
\mathcal I
\rightarrow
\mathcal U.
}
\]

即系统可以根据获得的信息产生控制：

\[
u_t
=
\Gamma(
x_t,
i_t
).
\]

于是形成最基本的运行闭环：

\[
\boxed{
Information
\rightarrow
StateEvaluation
\rightarrow
Control
\rightarrow
StateChange.
}
\]

控制信号的本质不在于它表达了什么语言内容，而在于：

\[
\boxed{
\text{它是否真正改变了系统的状态转移条件。}
}
\]

\section{信用信号：把结果重新分配给过去}

前三类信号已经足以描述一个能够运行的控制系统：

\[
State
+
Information
\rightarrow
Control
\rightarrow
NewState.
\]

但是还不能描述一个真正能够持续学习的 Agent。

原因是：

系统产生结果以后，还必须回答一个新的问题：

\[
\boxed{
\text{为什么产生这个结果？}
}
\]

进一步：

\[
\boxed{
\text{系统内部哪些结构应该因此发生变化？}
}
\]

这就引出了第四类信号：

\[
\boxed{
CreditSignal.
}
\]

这里的 Credit 不是经济意义上的信用，而是机器学习和控制意义上的：

\[
\boxed{
\text{责任归因}.
}
\]

设一次结果产生评价：

\[
r_t.
\]

可能表示：

成功；

失败；

预测残差；

资源代价；

约束违反；

长期收益。

仅仅知道：

\[
r_t=-1
\]

还不够。

系统必须进一步判断：

这个负结果究竟应当归因于：

错误表示？

错误策略？

错误记忆？

错误参数？

错误资源分配？

错误历史更新？

因此信用信号的基本定义为：

\[
\boxed{
\mathcal S_C
=
\text{把结果的学习责任分配给可塑性结构的信号}.
}
\]

设系统包含可塑性结构集合：

\[
\Xi
=
\{
\xi_1,\xi_2,\ldots,\xi_n
\}.
\]

则信用分配可以表示为：

\[
\boxed{
c_t
=
[
c_1(t),
c_2(t),
\ldots,
c_n(t)
].
}
\]

其中：

\[
c_i(t)
\]

表示第 $i$ 个结构针对当前结果应承担的学习责任。

\section{信用与误差不是同一个概念}

误差：

\[
\varepsilon_t
\]

描述：

\[
\boxed{
\text{发生了多大的偏差}.
}
\]

信用：

\[
c_i(t)
\]

描述：

\[
\boxed{
\text{这个偏差应该由谁承担}.
}
\]

因此：

\[
\boxed{
Error
\neq
Credit.
}
\]

例如：

\[
\varepsilon_t=10
\]

只表示预测出现很大错误。

但可能存在：

\[
c_1=0.8,
\qquad
c_2=0.1,
\qquad
c_3=0.1.
\]

表示这个错误主要应当修改：

\[
\xi_1.
\]

所以一般应有一个信用生成过程：

\[
\boxed{
c_t
=
R_C(
History_t,
Outcome_t,
\varepsilon_t
).
}
\]

其中：

\[
R_C
\]

是信用归因算子。

这意味着：

\[
\boxed{
\text{误差是结果描述，信用是学习路由。}
}
\]

\section{信用与控制也不是同一种信号}

控制信号回答：

\[
\boxed{
\text{现在应该怎样改变运行状态？}
}
\]

信用信号回答：

\[
\boxed{
\text{以后哪些结构应该因此发生学习？}
}
\]

因此：

\[
\boxed{
Control
\neq
Credit.
}
\]

可以用一个简单例子说明。

Agent 发现预测错误。

系统可能立即产生控制信号：

\[
u_t:
\text{增加搜索范围}.
\]

与此同时产生信用信号：

\[
c_t:
\text{某个长期结构可能需要重新学习}.
\]

前者改变：

\[
x_{t+\Delta t}.
\]

后者改变：

\[
\Xi_{t+\Delta T}.
\]

其中通常：

\[
\Delta T
\ge
\Delta t.
\]

因此控制主要属于：

\[
\boxed{
\text{运行动力学}}
\]

而信用主要属于：

\[
\boxed{
\text{学习动力学}}.
\]

\section{AgentOS 的双动力系统}

四类信号出现以后，可以把 AgentOS 分成两个彼此耦合的基本动力系统。

第一个是：

\[
\boxed{
\text{运行动力系统}.
}
\]

其基本形式为：

\[
\boxed{
x_{t+1}
=
F(
x_t,
i_t,
u_t;
\Xi_t
).
}
\]

其中：

状态：

\[
x_t
\]

信息：

\[
i_t
\]

控制：

\[
u_t
\]

共同决定下一状态。

第二个是：

\[
\boxed{
\text{学习动力系统}.
}
\]

其基本形式为：

\[
\boxed{
\Xi_{t+1}
=
G(
\Xi_t,
c_t
).
}
\]

其中：

\[
c_t
\]

决定哪些结构获得多少修改责任。

于是完整系统可以表示为：

\[
\boxed{
\begin{aligned}
x_{t+1}
&=
F(x_t,i_t,u_t;\Xi_t),
\\
\Xi_{t+1}
&=
G(\Xi_t,c_t).
\end{aligned}
}
\]

这是 AgentOS 一个非常基础的双动力方程。

\section{运行闭环与学习闭环}

由此可以进一步得到两个基本闭环。

运行闭环：

\[
\boxed{
State
\rightarrow
Information
\rightarrow
Control
\rightarrow
NewState.
}
\]

更准确地说：

\[
(x_t,i_t)
\rightarrow
u_t
\rightarrow
x_{t+1}.
\]

学习闭环：

\[
\boxed{
Outcome
\rightarrow
Credit
\rightarrow
StructureChange
\rightarrow
FutureBehaviorChange.
}
\]

即：

\[
Outcome_t
\rightarrow
c_t
\rightarrow
\Xi_{t+1}
\rightarrow
F_{\Xi_{t+1}}.
\]

所以 AgentOS 并不是一个单闭环系统，而是：

\[
\boxed{
\text{运行闭环}
\bowtie
\text{学习闭环}.
}
\]

运行闭环解决：

\[
\boxed{
\text{现在怎样行动}.
}
\]

学习闭环解决：

\[
\boxed{
\text{以后怎样变得不同}.
}
\]

\section{四类信号之间的基本转换关系}

四类信号并不是彼此封闭的。

它们之间存在稳定的转换关系。

首先：

\[
\boxed{
Information
\rightarrow
State.
}
\]

系统获得新信息以后：

\[
i_t
\]

会改变当前状态：

\[
x_t.
\]

其次：

\[
\boxed{
State+Information
\rightarrow
Control.
}
\]

控制一般可以写成：

\[
u_t
=
\Pi(
x_t,
i_t
).
\]

第三：

\[
\boxed{
StateTransition
\rightarrow
Outcome.
}
\]

即：

\[
(x_t,u_t,x_{t+1})
\rightarrow
r_t.
\]

第四：

\[
\boxed{
Outcome+History
\rightarrow
Credit.
}
\]

即：

\[
c_t
=
R_C(
x_{0:t},
i_{0:t},
u_{0:t},
r_t
).
\]

第五：

\[
\boxed{
Credit
\rightarrow
StructureChange.
}
\]

即：

\[
\Delta\Xi_t
=
A(
\Xi_t,
c_t
).
\]

最终结构改变又影响下一轮运行：

\[
\boxed{
\Xi_{t+1}
\rightarrow
F_{t+1}.
}
\]

因此完整循环为：

\[
\boxed{
Information
\rightarrow
State
\rightarrow
Control
\rightarrow
Outcome
\rightarrow
Credit
\rightarrow
Structure
\rightarrow
NewStateDynamics.
}
\]

这构成 AgentOS 最基本的长期循环之一。

\section{四类信号不是严格一一对应于模块}

建立信号本体以后，必须避免另一个错误：

\[
\boxed{
\text{一种信号}
=
\text{一个模块}.
}
\]

真实 AgentOS 中，一个模块可能：

接收状态；

读取信息；

产生控制；

同时产生信用相关证据。

因此更合理的模块定义是：

\[
\boxed{
M:
(\mathcal S_S,
\mathcal S_I,
\mathcal S_U,
\mathcal S_C)
\rightarrow
(\mathcal S_S',
\mathcal S_I',
\mathcal S_U',
\mathcal S_C').
}
\]

模块只是不同信号之间的变换节点。

所以 AgentOS 的整体结构不应该只理解成：

\[
ModuleGraph.
\]

还应理解成：

\[
\boxed{
SignalFlowGraph.
}
\]

模块图描述：

\[
\text{谁连接谁}.
\]

信号图描述：

\[
\boxed{
\text{什么性质的因果量正在流动}.
}
\]

后者对于理解智能运行时更加基础。

\section{混合信号与主导因果角色}

现实系统中，一个具体消息可能同时具有多种作用。

例如外部反馈：

\[
\text{“任务失败，因为目标对象位置判断错误。”}
\]

它首先包含信息：

\[
i_t.
\]

它可能改变当前状态：

\[
x_t.
\]

它也可能触发控制调整：

\[
u_t.
\]

同时还可能成为信用生成依据：

\[
c_t.
\]

因此，四类信号不是要求每个物理消息只能拥有一个类别。

更准确的描述是：

\[
\boxed{
\sigma_t
=
[
w_S,
w_I,
w_U,
w_C
]
}
\]

其中：

\[
w_S+w_I+w_U+w_C=1
\]

可以表示一个复合信号在四种因果角色上的权重。

但工程实现中，更推荐把不同作用显式拆开：

\[
\sigma
\rightarrow
\{
\sigma_S,
\sigma_I,
\sigma_U,
\sigma_C
\}.
\]

这样可以避免把：

内容；

状态；

命令；

学习责任

混成一条不可治理的数据流。

\section{时间尺度是四类信号之外的第二维度}

信号类型与时间尺度必须正交。

不能简单认为：

状态一定快；

信用一定慢。

虽然通常存在这样的倾向，但并非绝对。

定义：

\[
\tau(\sigma)
\]

为信号有效时间尺度。

于是任何信号都具有：

\[
\boxed{
Type
\times
Timescale
}
\]

两个属性。

例如：

快速状态信号；

慢状态信号；

即时控制；

长期控制；

即时信用；

延迟信用。

因此 AgentOS 的信号结构至少是二维的：

\[
\boxed{
SignalOntology
\times
TemporalScale.
}
\]

进一步结合功能拓扑，可以形成：

\[
\boxed{
AgentOS
=
FunctionalTopology
\times
TemporalHierarchy
\times
SignalOntology.
}
\]

这三个维度分别描述：

\[
\boxed{
\text{在哪里发生}
}
\]

\[
\boxed{
\text{以什么时间尺度发生}
}
\]

\[
\boxed{
\text{以什么因果角色发生}.
}
\]

\section{信号边界与信号权限}

并不是任何信号都应该拥有相同的系统权限。

例如，一个信息信号可以被读取，并不意味着它能够直接变成：

控制；

信用；

长期结构改变。

因此应当区分：

\[
P_I
\]

信息进入权限；

\[
P_U
\]

控制作用权限；

\[
P_C
\]

信用产生权限；

\[
P_L
\]

结构修改权限。

一般不存在：

\[
P_I
\Rightarrow
P_U.
\]

也不存在：

\[
P_I
\Rightarrow
P_C.
\]

更不存在：

\[
P_I
\Rightarrow
P_L.
\]

因此可以形成一条基本安全原则：

\[
\boxed{
\textbf{
能够被认知的内容，不自动拥有控制系统或修改系统结构的权力。
}
}
\]

这是 AgentOS 区别于简单信息处理流水线的重要原则。

\section{信用信号为什么必须独立成为第四类信号}

可以进一步问：

为什么不能把信用看成普通信息？

因为：

\[
c_t
\]

当然也可以被编码成数字或向量。

但如果只把它理解为信息，就无法表达其特殊因果权力。

信用信号的特殊之处在于：

\[
\boxed{
\text{它具有结构修改指向性}.
}
\]

即：

\[
c_t
\rightarrow
\Delta\Xi.
\]

而普通信息：

\[
i_t
\]

默认只参与：

\[
x_{t+1}.
\]

所以：

\[
\boxed{
Information
\rightarrow
RuntimeInterpretation
}
\]

而：

\[
\boxed{
Credit
\rightarrow
PlasticityAllocation.
}
\]

正因为持续智能不仅需要运行，还需要决定：

\[
\boxed{
\text{哪些结构应该因现实反馈而改变},
}
\]

信用才必须从信息中独立出来，成为 AgentOS 的第四类基本信号。

\section{信用的局部性与延迟性}

信用信号具有两个与其他信号不同的重要性质。

第一是延迟性。

某个结构在：

\[
t_0
\]

发生作用，

真正结果可能在：

\[
t_1\gg t_0
\]

才出现。

所以：

\[
\boxed{
CauseTime
\neq
CreditTime.
}
\]

第二是局部性。

一个全局失败：

\[
Failure
\]

不意味着所有内部结构都应该修改。

因此：

\[
\boxed{
GlobalError
\neq
GlobalPlasticity.
}
\]

合理系统必须找到：

\[
c_i
\]

使：

\[
\Delta\xi_i
\propto
c_i.
\]

因此信用信号天然要求：

历史；

责任；

路径；

依赖

共同参与计算。

这也是持续智能体学习问题比普通即时控制更困难的原因之一。

\section{四类信号与因果方向}

四类信号还可以按照主要因果方向进一步理解。

状态信号：

\[
\boxed{
\text{横截面描述}
}
\]

描述当前系统。

信息信号：

\[
\boxed{
\text{表征输入}
}
\]

把外部或内部结构带入当前计算。

控制信号：

\[
\boxed{
\text{前向因果}
}
\]

影响未来状态。

信用信号：

\[
\boxed{
\text{反向因果归因}
}
\]

根据未来结果重新评价过去结构。

这里的“反向”不是物理时间倒流。

真实物理过程仍然是：

\[
t\rightarrow t+\Delta t.
\]

所谓反向，是计算意义上的：

\[
\boxed{
FutureOutcome
\rightarrow
PastResponsibilityAssignment.
}
\]

因此可以把四种角色压缩成：

\[
\boxed{
State:
Being
}
\]

\[
\boxed{
Information:
Representation
}
\]

\[
\boxed{
Control:
ForwardInfluence
}
\]

\[
\boxed{
Credit:
BackwardAttribution
}
\]

这四者共同形成一个完整的持续智能运行环。

\section{最小 AgentOS 信号系统}

一个最小持续 AgentOS 至少需要：

当前状态：

\[
x_t;
\]

输入信息：

\[
i_t;
\]

控制策略：

\[
u_t
=
\Pi(x_t,i_t);
\]

状态动力学：

\[
x_{t+1}
=
F(x_t,i_t,u_t;\Xi_t);
\]

结果评价：

\[
r_t
=
R(
x_t,
u_t,
x_{t+1}
);
\]

信用生成：

\[
c_t
=
C(
History_t,
r_t
);
\]

结构适应：

\[
\Xi_{t+1}
=
G(
\Xi_t,
c_t
).
\]

因此最小闭环可以写成：

\[
\boxed{
\begin{aligned}
i_t
&\rightarrow x_t,
\\
(x_t,i_t)
&\rightarrow u_t,
\\
(x_t,i_t,u_t)
&\rightarrow x_{t+1},
\\
(x_t,u_t,x_{t+1})
&\rightarrow r_t,
\\
(History_t,r_t)
&\rightarrow c_t,
\\
(\Xi_t,c_t)
&\rightarrow\Xi_{t+1}.
\end{aligned}
}
\]

这已经具备：

感知；

状态；

控制；

结果；

归因；

适应

六个最小环节。

\section{从四类信号重新理解智能}

四信号框架还提供了一种重新理解智能的方法。

只有状态而没有信息：

系统能够存在，但无法形成丰富表示。

只有状态和信息：

系统能够表示，却不一定能够主动改变运行方向。

加入控制以后：

系统能够：

\[
\boxed{
\text{根据表示改变未来.}
}
\]

但仍然可能永远以相同方式行动。

只有再加入信用：

系统才能：

\[
\boxed{
\text{根据行动结果改变以后如何行动.}
}
\]

因此可以得到一个递进关系：

\[
\boxed{
State
}
\]

提供存在基础；

\[
\boxed{
State+Information
}
\]

产生表征能力；

\[
\boxed{
State+Information+Control
}
\]

产生受表示驱动的主动作用；

\[
\boxed{
State+Information+Control+Credit
}
\]

产生持续适应能力。

所以对于持续智能：

\[
\boxed{
\textbf{
信用不是附加信号，而是把“会行动的系统”变成“会因行动结果而改变自己的系统”的关键闭环。
}
}
\]

\section{AgentOS 信号本体的第一性表述}

至此，可以把 AgentOS 运行时的四种基本因果角色统一表示为：

\[
\boxed{
\mathfrak S
=
(
X,I,U,C
)
}
\]

其中：

\[
X
\]

为状态空间；

\[
I
\]

为信息空间；

\[
U
\]

为控制空间；

\[
C
\]

为信用空间。

整个 AgentOS 可以抽象为两个耦合映射：

运行映射：

\[
\boxed{
F:
X\times I\times U\times\Xi
\rightarrow X
}
\]

和学习映射：

\[
\boxed{
G:
\Xi\times C
\rightarrow\Xi.
}
\]

控制由：

\[
\boxed{
\Pi:
X\times I
\rightarrow U
}
\]

产生；

信用由：

\[
\boxed{
R_C:
History\times Outcome
\rightarrow C
}
\]

产生。

因此完整形式为：

\[
\boxed{
\begin{aligned}
u_t
&=
\Pi(x_t,i_t),
\\
x_{t+1}
&=
F(x_t,i_t,u_t;\Xi_t),
\\
c_t
&=
R_C(H_t,o_{t+1}),
\\
\Xi_{t+1}
&=
G(\Xi_t,c_t).
\end{aligned}
}
\]

这四个方程构成 AgentOS 最小的运行--学习统一动力结构。

\section{本章结论：AgentOS 是四类信号的多时间尺度耦合系统}

传统计算系统主要关注：

\[
Data
\rightarrow
Computation
\rightarrow
Output.
\]

持续智能体则必须同时处理：

\[
\boxed{
State,
Information,
Control,
Credit.
}
\]

其中：

\[
\boxed{
State
}
\]

回答系统现在是什么；

\[
\boxed{
Information
}
\]

回答系统当前表示了什么；

\[
\boxed{
Control
}
\]

回答系统接下来怎样改变；

\[
\boxed{
Credit
}
\]

回答哪些过去结构应该因为结果而改变。

四者形成：

\[
\boxed{
(State,Information)
\rightarrow
Control
\rightarrow
State'
\rightarrow
Outcome
\rightarrow
Credit
\rightarrow
Structure'
}
\]

而：

\[
Structure'
\]

进一步改变下一轮状态动力学：

\[
\boxed{
Structure'
\rightarrow
(State',Information')
\rightarrow
Control'.
}
\]

因此可以给 AgentOS 一个更基础的运行时定义：

\[
\boxed{
\textbf{
AgentOS 是状态、信息、控制与信用四类信号
在多时间尺度功能结构中持续传播、转换和闭合的智能运行时。
}
}
\]

进一步：

\[
\boxed{
\textbf{
状态承担存在，
信息承担表征，
控制承担作用，
信用承担学习。
}
}
\]

这四种因果角色共同构成了持续智能体从：

\[
\text{存在}
\rightarrow
\text{认识}
\rightarrow
\text{行动}
\rightarrow
\text{改变自己}
\]

的最小完整闭环。

因此，后续 AgentOS 的任何具体结构设计，都应能够回答四个问题：

\[
\boxed{
\text{它读取和改变哪些状态？}
}
\]

\[
\boxed{
\text{它接收和产生哪些信息？}
}
\]

\[
\boxed{
\text{它能够产生哪些控制作用？}
}
\]

\[
\boxed{
\text{它如何产生、接收或传递学习信用？}
}
\]

当一个模块无法回答这四个问题时，它在 AgentOS 中的动力学位置仍然是不完整的。
